\documentclass{article}
\usepackage{amsmath}
\usepackage{amssymb}

\begin{document}

\textbf{Problem 2.} Given positive integers $k, m, n, p$ such that $p = 2^{2^n} + 1$, $p$ is a prime number, and $2^k - m$ is divisible by $p$. Prove that there exists a positive integer $\ell$ such that the number $2^\ell - m$ is divisible by $p^2$.

Here is the step-by-step solution to the problem.

### Proof

Let $p = 2^{2^n} + 1$ be a Fermat prime with $n \in \mathbb{Z}^+$.
We are given that $2^k \equiv m \pmod{p}$.
This implies that there exists an integer $s$ such that:
$$ m = 2^k + s \cdot p $$

We want to find a positive integer $\ell$ such that $2^\ell \equiv m \pmod{p^2}$.

**Step 1: Determine the behavior of powers of 2 modulo $p^2$.**

First, observe the order of $2$ modulo $p$.
Since $p = 2^{2^n} + 1$, we have:
$$ 2^{2^n} = p - 1 \equiv -1 \pmod{p} $$
Squaring both sides:
$$ 2^{2^{n+1}} \equiv (-1)^2 \equiv 1 \pmod{p} $$
Let $\delta = 2^{n+1}$. This is the order of $2$ modulo $p$.
Now, let's examine $2^\delta$ modulo $p^2$:
$$ 2^\delta = 2^{2^{n+1}} = (2^{2^n})^2 = (p - 1)^2 $$
Expanding the square:
$$ (p - 1)^2 = p^2 - 2p + 1 $$
Considering this modulo $p^2$:
$$ 2^\delta \equiv -2p + 1 \pmod{p^2} $$
$$ 2^\delta = 1 - 2p \pmod{p^2} $$

**Step 2: Construct $\ell$.**

We look for $\ell$ in the form $\ell = k + x \cdot \delta$, where $x$ is a positive integer to be determined.
$$ 2^\ell = 2^{k + x\delta} = 2^k \cdot (2^\delta)^x $$
Substituting the equivalence found in Step 1:
$$ 2^\ell \equiv 2^k \cdot (1 - 2p)^x \pmod{p^2} $$
Using the Binomial Theorem, $(1 - 2p)^x = 1 - x(2p) + \binom{x}{2}(2p)^2 - \dots$.
Modulo $p^2$, all terms containing $p^2$ or higher powers vanish. Thus:
$$ (1 - 2p)^x \equiv 1 - 2xp \pmod{p^2} $$
Substituting this back into the expression for $2^\ell$:
$$ 2^\ell \equiv 2^k (1 - 2xp) \pmod{p^2} $$
$$ 2^\ell \equiv 2^k - x \cdot 2^{k+1}p \pmod{p^2} $$

**Step 3: Solve for $x$.**

We require $2^\ell \equiv m \pmod{p^2}$.
Recall that $m = 2^k + sp$. Therefore, we need:
$$ 2^k - x \cdot 2^{k+1}p \equiv 2^k + sp \pmod{p^2} $$
Subtracting $2^k$ from both sides:
$$ -x \cdot 2^{k+1}p \equiv sp \pmod{p^2} $$
Since $p \neq 0$, we can divide the congruence by $p$ (which changes the modulus from $p^2$ to $p$):
$$ -x \cdot 2^{k+1} \equiv s \pmod{p} $$
We need to solve this linear congruence for $x$.
Since $p$ is an odd prime (as $p = 2^{2^n} + 1$), $\gcd(2^{k+1}, p) = 1$.
Therefore, $2^{k+1}$ has a multiplicative inverse modulo $p$.
Let $A = -2^{k+1}$. The congruence is $A \cdot x \equiv s \pmod{p}$.
There exists a unique solution for $x$ modulo $p$:
$$ x \equiv s \cdot A^{-1} \pmod{p} $$
We can choose a positive integer representative for $x$. Let this specific value be $x_0$.

**Conclusion**

By choosing $\ell = k + x_0 \cdot 2^{n+1}$, we satisfy the condition:
$$ 2^\ell - m \equiv 0 \pmod{p^2} $$
Since $k$ and $x_0$ are positive, $\ell$ is a positive integer.
Thus, such a number $\ell$ exists.

Q.E.D.

\end{document}